% Independent contribution. AI-produced by GPT-6 Astra; no priority claimed.
\subsection{The exact finite lifting problem for the literal fixed subgroup}

We retain $G=\widehat{\mathrm{GT}}_0$, its multiplier-one
subgroup $H$, and the literal continuous endomorphism $D=D_P$.
The normalized peripheral action identifies these compact groups
with their faithful outer actions on $B=\widehat{C_2*C_3}$.
All finite quotients below are quotients of this full profinite
group.  Let $q_m:B\twoheadrightarrow Q_m$ be a nested cofinal
sequence of finite characteristic quotients retaining both
$B/F=\mathfrak S_3$ and $B^{\rm ab}=C_6$.
Write
\[
 \rho_m:G\longrightarrow\operatorname{Out}(Q_m)
 \tag{B.X1}
\]
for the induced continuous homomorphism.  Enlarge the initial
quotient, if necessary, so that the continuous character
$\Delta|_H$ factors through $\rho_1|_H$; denote the resulting
character of $\rho_m(H)$ by $\delta_m$.
This choice is possible because the characteristic finite
actions give the topology of the faithful compact outer image:
an outer action trivial at every finite level is inner, by
compactness of the finite conjugator conditions.

Define the actual finite joint source
\[
 \begin{gathered}
 K_m=\ker q_m\cap B'\cap
       \operatorname{core}_{B}\ker(q_m\varphi_P),
       \qquad R_m=B/K_m,\\
 R_m^c=B/\overline{[B,K_m]},\qquad
 M_m=K_m/\overline{[B,K_m]} .
 \end{gathered}                                      \tag{B.X2}
\]
The core is taken from the named index-six subgroup $F$ of $B$.
Thus $K_m$ is open and normal in $B$, and $R_m$ has actual
maps to $Q_m$ both by reduction and, on its pure subgroup,
by the descended map $q_m\varphi_P$.  The finite-cover proof
of B.2.7 makes $R_m^c$ finite.  Let $\mathcal A_m$ be its
finite group of positive, pure-framed, cusp-fixed automorphisms
preserving $M_m$, in the exact sense of B.2.9.
Here an element includes its automorphism and its two
pure torsion conjugators, with the multiplication (B.Q3).
There are two homomorphisms on this whole group,
\[
 p_m,d_m:\mathcal A_m\longrightarrow\operatorname{Out}(Q_m).
 \tag{B.X3}
\]
Here $p_m$ is ordinary reduction.  The second map first
reduces to $R_m$, restricts to its pure subgroup, and uses
$q_m\varphi_P$ to descend, with the normalized target inner
ambiguity removed by passing to the displayed outer group.
These operations are defined on all of $\mathcal A_m$:
B.2.7 gives an actual full positive source lift, and
characteristicness of $\ker q_m$ then proves preservation
of both required kernels.  The two resulting outer classes
depend only on the retained finite automorphism.  In particular,
\[
 (p_m,d_m)(\mathcal A_m)
   =\{(\rho_m(h),\rho_m(Dh)):h\in H\}.
 \tag{B.X4}
\]
This is an equality of actual finite images, not merely a
test of the two cycle equations.

Put $\mathcal G_m=\rho_m(G_{\mathbb Q})$ and $c_m=\rho_m(c)$.
The arithmetic subgroup in this definition is the fixed
Drinfeld--Ihara image.  For $v\in\mathcal G_m$ put
$b_m(v)=vc_mv^{-1}c_m$, and define two explicitly finite
subsets of $\mathbb F_2$:
\[
\begin{aligned}
 S_m^H&=\{\delta_m(p_m(a)):a\in\mathcal A_m,
                                  \ d_m(a)=p_m(a)\},\\
 S_m^G&=\{\delta_m(p_m(a)):
       a\in\mathcal A_m,\ v\in\mathcal G_m,
                    \ d_m(a)=p_m(a)b_m(v)\}.
\end{aligned}                                        \tag{B.X5}
\]
The values $p_m(a)$ lie in $\rho_m(H)$ by (B.X4), so
every occurrence of $\delta_m$ is typed on its actual domain.

\begin{result}{B.1.9}{Proved}{Exact finite images controlling fixed-point parity}
For the literal fixed subgroup $C_P=\operatorname{Fix}(D)$,
\[
 \Delta(C_P\cap H)=\bigcap_{m\ge1}S_m^H,
 \qquad
 \Delta(C_P)=\bigcap_{m\ge1}S_m^G.
 \tag{B.X6}
\]
The sets on each right-hand side decrease with $m$.
Thus the full-multiplier problem retains the actual
arithmetic commutator $b_m(v)$; it does not replace it
by the multiplier-one equalizer.
\end{result}
\begin{proof}
For the first equality consider the closed subsets
\[
 X_m(e)=\{h\in H:\Delta(h)=e,
                   \rho_m(Dh)=\rho_m(h)\}
 \tag{B.X7}
\]
of the compact group $H$.  Equation (B.X4) proves that
their nonemptiness is exactly $e\in S_m^H$.
They decrease with $m$.  Their intersection consists
precisely of the elements of $C_P\cap H$ of parity $e$,
because the faithful finite outer actions separate
elements of $G$.  Compactness proves the equality.

For the full group use the actual surjection of compact
spaces $H\times G_{\mathbb Q}\to G$, $(h,\gamma)\mapsto
h\gamma$: the arithmetic cyclotomic character is onto
$\widehat{\mathbb Z}^{\times}$.  Arithmetic descent is
$D\gamma=c\gamma c$, so multiplication in $G$ gives the
exact identity
\[
 D(h\gamma)=h\gamma
 \quad\Longleftrightarrow\quad
 Dh=h b_\gamma,
 \qquad b_\gamma=\gamma c\gamma^{-1}c\in H\cap G_{\mathbb Q}.
 \tag{B.X8}
\]
Also $\Delta(h\gamma)=\Delta(h)$, by arithmetic vanishing.
The closed sets in $H\times G_{\mathbb Q}$ defined by
$\Delta(h)=e$ and
$\rho_m(Dh)=\rho_m(h)\rho_m(b_\gamma)$ are nonempty
exactly for $e\in S_m^G$: (B.X4) lifts $a$ to an actual
$h$, and the definition of $\mathcal G_m$ lifts $v$ to
an actual $\gamma$, independently.  They are nested.
Their intersection is precisely the inverse image of
the parity-$e$ part of $C_P$ under the displayed
surjection.  Compactness proves the second equality.
Every construction commutes with a finer characteristic
quotient, proving the nesting assertion.
\end{proof}

The full arithmetic correction also has the useful exact
identities
\[
 D(b_\gamma)=b_\gamma^{-1},\qquad
 Dh=h b_\gamma\ \Longrightarrow\ D^2h=h.
 \tag{B.X9}
\]
They follow by substituting $D\gamma=c\gamma c$ and
$Dc=c$.  The implication in (B.X9) is only a necessary
consequence; (B.X8), rather than the weaker periodicity
equation, is the full fixed-point condition.

\paragraph{A restriction on a possible character-coboundary proof.}
The actual element $k$ of B.1.3 can be chosen so that its
normalized action is the identity on every maximal pro-$p$
quotient $F^{(p)}$ of $F$, while $Dk=1$ and $\Delta(k)=1$.
Indeed $f_k\in R_{\rm nil}(F'F^3)$.  For every finite
$p$-group quotient $F\to P$, the image of $F'F^3$ is
nilpotent, so the image of $f_k$ is one.  Passing to the
inverse limit gives $f_k=1$ in $F^{(p)}$; the normalized
formulas fix $x$ and conjugate $y$ by this image.  Therefore
every continuous character $\eta:H\to\mathbb F_2$ factoring
through the product of these actual pro-$p$ automorphism
images satisfies
\[
                 (\eta+\eta\circ D)(k)=0\ne\Delta(k).
 \tag{B.X10}
\]
Consequently none of these characters supplies the identity
$\Delta=\eta+\eta\circ D$ on $H$.  This excludes precisely
the displayed comparison maps; it makes no assertion about
characters detected by other mixed finite quotients.

The exact sets (B.X5) have not all been computed.  In
particular the surjectivity of $D$ and of the character
map in B.2.25 has not been substituted for the remaining
finite equalizer tests in (B.X6).

For an explicit check of that logical limitation, put
$A=\prod_{n\ge0}\mathbb F_2$ and let $T:A\to A$ be
the left shift.  Use the two continuous characters
$e(a)=a_0$ and $e'(a)=a_0+a_1$.  In both cases
$(T,e)$, respectively $(T,e')$, is onto $A\times\mathbb F_2$:
for a prescribed tail $b$ and bit $v$, choose the missing
coordinate as $v$, respectively $v+b_0$.
The iterates are independent: a finite relation among
$a_n$ has its last nonzero coefficient unchanged, while
a relation among $a_n+a_{n+1}$ has the same nonzero
coefficient at its last appearing coordinate.  The
fixed elements of $T$ are exactly the constant sequences.
Thus $e(\operatorname{Fix}T)=\mathbb F_2$ and
$e'(\operatorname{Fix}T)=0$.  These statements concern
the two specified continuous morphisms of the explicit
compact group $A$.  The executable check
\path{check_shift_fixed_parity.py} verifies their finite
prefix maps and ranks through nine coordinates; the
coordinate argument just given proves the infinite
assertions.  No map from these examples to $G$ is asserted.
